\chapter{Memoria Descriptiva}

\section{Descripción general del proyecto}

El presente proyecto corresponde a la implementación de un sistema de cableado estructurado, red de datos, CCTV y conectividad inalámbrica para los \totalPisos{} niveles del pabellón de la Escuela Profesional de Ingeniería Informática y de Sistemas de la Universidad Nacional de San Antonio Abad del Cusco, ubicada en la Ciudad Universitaria de Perayoc. La memoria se desarrolla tomando como base los planos disponibles de la Escuela Profesional, los cuales presentan la estructura del pabellón y sirven como soporte para el trazado de canalizaciones, rutas y puntos de red.

El proyecto contempla la instalación de un backbone de fibra óptica multimodo OM4 que interconecta el Distribuidor Principal (MDF) ubicado en el segundo piso con los Distribuidores de Piso (IDF) del primer, tercer y cuarto nivel, cableado horizontal de cobre F/UTP Categoría 6A, puntos de telecomunicaciones en ambientes académicos y administrativos, gabinetes de comunicaciones, switches de acceso PoE+, así como puntos de acceso inalámbrico WiFi 6 para cobertura. Se consideran asimismo las rutas de escalerillas tipo malla, canaletas PVC, tuberías conduit y pasos técnicos para asegurar un despliegue ordenado y compatible con la infraestructura existente.

La intervención no modifica la estructura arquitectónica, sino que incorpora infraestructura tecnológica para asegurar conectividad, disponibilidad y crecimiento futuro. Se propone una implementación por niveles, definiendo claramente rutas, áreas técnicas y puntos de servicio en cada piso. El proyecto contempla un total consolidado de \textbf{\totalPuntos{} puntos de red}, garantizando plena consistencia de principio a fin del documento.

\section{Objetivos de la obra}

El objetivo general de la obra es diseñar e implementar un sistema de cableado estructurado y red de datos para la Facultad de Ingeniería Eléctrica, Electrónica, Informática y Mecánica de la UNSAAC, garantizando conectividad confiable, escalabilidad y cumplimiento de estándares técnicos.

Los objetivos específicos son los siguientes:

\begin{itemize}
    \item Levantar información de ambientes y necesidades de conectividad para definir la cantidad y ubicación de puntos de red y puntos de acceso inalámbrico.
    \item Diseñar la topología de red (backbone y cableado horizontal), rutas de canalización y ubicación de gabinetes de telecomunicaciones.
    \item Instalar cableado de cobre y fibra óptica, con terminaciones certificadas y ordenadas en patch panels y tomas de usuario.
    \item Implementar infraestructura activa de red (switches y puntos de acceso) para brindar conectividad cableada e inalámbrica.
    \item Aplicar criterios de administración, etiquetado y documentación para facilitar la operación y el mantenimiento.
    \item Garantizar compatibilidad con la infraestructura existente sin afectar la estructura arquitectónica.
\end{itemize}

\section{Ubicación y área de intervención}

El proyecto se ubica en el pabellón de la Escuela Profesional de Ingeniería Informática y de Sistemas de la Universidad Nacional de San Antonio Abad del Cusco, dentro de la Ciudad Universitaria de Perayoc, distrito, provincia y departamento del Cusco.

El área de intervención comprende los \totalPisos{} niveles del pabellón de la Escuela Profesional de Ingeniería Informática y de Sistemas, abarcando los ambientes académicos, administrativos y de servicios descritos a continuación:

\begin{itemize}
    \item \textbf{Primer nivel (Aulas y áreas administrativas):} alberga un Distribuidor de Piso (\textbf{IDF-01} de 24U) secundario ubicado debajo de la escalera principal (espacio acondicionado). Los ambientes intervenidos incluyen las aulas académicas principales, la oficina de la Secretaría de Escuela y la oficina de la Dirección de Escuela. Se proyectan \textbf{20 puntos de red} en total (incluyendo 16 puertos para datos correspondientes a 8 faceplates dobles y 4 puntos para APs de cobertura inalámbrica, con 0 cámaras de videovigilancia CCTV).
    \item \textbf{Segundo nivel (Nivel principal de red):} alberga el Distribuidor Principal de Cableado (\textbf{MDF-02} de 42U) del proyecto, desde donde se distribuye el backbone hacia los pisos superior e inferior, y 2 gabinetes de zona (ZG). Los ambientes intervenidos incluyen aulas académicas, biblioteca y laboratorio. Se proyectan \textbf{85 puntos} de conexión en total, equivalentes a \textbf{85 puertos/cables de red reales} (incluyendo 78 puertos correspondientes a 39 faceplates dobles, 3 puntos para APs de cobertura inalámbrica y 4 para cámaras de videovigilancia CCTV).
    \item \textbf{Tercer nivel (Alta densidad académica):} concentra laboratorios de cómputo de la escuela (ambientes 301 a 306 y 309 a 312), la sala de reunión de docentes, la sala de exposiciones y la jefatura del Centro de Capacitación en Informática (CCI). Cuenta con un Distribuidor de Piso (\textbf{IDF-03} de 24U) y \textbf{10 gabinetes de zona (ZG)} (uno por laboratorio). Se proyectan \textbf{448 puntos} en total, equivalentes a \textbf{448 puertos/cables de red reales} (incluyendo 448 puertos correspondientes a 224 faceplates dobles, con 0 APs y 0 cámaras de videovigilancia CCTV tras la reubicación de puntos a pared).
    \item \textbf{Cuarto nivel (Gestión e investigación):} comprende los cubículos de docentes, el auditorio principal, la cabina de control de cámaras, el laboratorio de investigación, los círculos de estudio, la secretaría de escuela y la dirección de la escuela. Cuenta con un Distribuidor de Piso (\textbf{IDF-04} de 24U) y \textbf{2 gabinetes de zona (ZG)}. Se proyectan \textbf{128 puntos} de red en total (incluyendo 120 puertos correspondientes a 60 faceplates dobles, 0 APs y 8 para cámaras de videovigilancia CCTV).
\end{itemize}

El alcance total de puntos de red reales (enlaces horizontales a certificar) asciende a \textbf{\totalPuntos{} puntos} distribuidos en los \totalPisos{} niveles, correspondientes a 331 tomas dobles de datos y 19 dispositivos de servicio (AP y CCTV). La intervención se concentra en el trazado de rutas de cableado, ubicación de gabinetes y puntos de red, sin alterar la estructura arquitectónica existente.

\section{Alcances del proyecto}

El alcance del proyecto comprende la implementación de la infraestructura de cableado y red en los \totalPisos{} niveles del pabellón de la Escuela Profesional de Ingeniería Informática y de Sistemas, considerando la instalación física, certificación y configuración básica para operación.

La intervención incluye:

\begin{itemize}
    \item Levantamiento de información técnica y definición de \textbf{\totalPuntos{} puntos de red reales} (enlaces) y las ubicaciones físicas distribuidas en los \totalPisos{} niveles de intervención.
    \item Diseño de rutas de canalización principal (escalerillas tipo malla), canalización secundaria (canaletas PVC con división física), tuberías conduit para backbone y registros de paso.
    \item Suministro e instalación de cableado horizontal \textbf{F/UTP Cat 6A LSZH} para datos, telefonía IP y CCTV.
    \item Suministro e instalación de cableado backbone de \textbf{fibra óptica multimodo OM4} entre el MDF (segundo piso) y los IDF (primer, tercer y cuarto piso).
    \item Suministro e instalación de \textbf{cuatro gabinetes principales} de telecomunicaciones: MDF-02 de 42U en el segundo piso, IDF-01 de 24U en el primer piso, IDF-03 de 24U en el tercer piso e IDF-04 de 24U en el cuarto piso, y \textbf{14 gabinetes de zona (ZG)}, incluyendo patch panels, organizadores y sistema de puesta a tierra.
    \item Instalación de tomas de usuario (faceplates de 2 puertos con iconografía, jacks Cat 6A blindados, cajas superficiales PVC) y su respectivo etiquetado bajo norma ANSI/TIA-606-C.
    \item Implementación de equipos activos de red (switches PoE+ Cisco Catalyst 9200, puntos de acceso WiFi 6, cámaras IP de vigilancia, NVR y UPS).
    \item Sistema completo de \textbf{puesta a tierra para telecomunicaciones} (TMGB, TGB, conductores 6 AWG y bonding de estructuras metálicas), conforme a ANSI/TIA-607-D.
    \item Pruebas, certificación del 100\% de los enlaces con equipo Fluke DSX-8000 y documentación de la red instalada.
    \item Elaboración de planos As-Built, planos de racks (front-view), diagramas de backbone, detalles constructivos de canalizaciones y diagrama de puesta a tierra.
    \item Entrega de dossier de certificación, manual de operación y mantenimiento, y documentación de administración del cableado.
\end{itemize}

El proyecto no contempla modificaciones estructurales mayores del edificio; las intervenciones se limitan a canalizaciones adosadas, soportes y pasos técnicos necesarios para el cableado.

\section{Justificación técnica y beneficios}

La justificación técnica del proyecto se basa en la necesidad de dotar a la facultad de una infraestructura de red confiable y estandarizada que soporte los servicios académicos, administrativos y de investigación. La conectividad es un componente esencial para el funcionamiento de laboratorios, aulas digitales, plataformas institucionales y servicios de comunicaciones.

Desde el punto de vista técnico, la intervención se justifica por los siguientes aspectos:

\begin{itemize}
    \item Necesidad de garantizar conectividad estable y de alta disponibilidad en todos los ambientes.
    \item Requerimiento de un sistema de cableado estructurado que permita crecimiento y mantenimiento ordenado.
    \item Integración de red cableada e inalámbrica con cobertura adecuada para usuarios móviles.
    \item Implementación de estándares de calidad que aseguren desempeño y compatibilidad futura.
    \item Reducción de fallas operativas mediante administración y etiquetado correcto del cableado.
\end{itemize}

Los beneficios esperados son los siguientes:

\begin{itemize}
    \item Mejora del acceso a servicios digitales para estudiantes, docentes y personal administrativo.
    \item Mayor confiabilidad en la transmisión de datos y continuidad de servicios.
    \item Facilidad de mantenimiento y escalabilidad de la red institucional.
    \item Ordenamiento de infraestructura tecnológica con documentación actualizada.
    \item Soporte para futuras ampliaciones y modernización de equipos de red.
\end{itemize}

\section{Normas y estándares aplicados}

La ejecución del proyecto debe desarrollarse considerando normas y criterios técnicos vinculados a cableado estructurado, telecomunicaciones, seguridad eléctrica y administración de infraestructura. Para ello, se toman como referencia las siguientes normas y documentos técnicos:

\begin{itemize}
    \item \textbf{ANSI/TIA-568.0-D:} cableado genérico para instalaciones de telecomunicaciones.
    \item \textbf{ANSI/TIA-568.1-D:} estándar de cableado comercial.
    \item \textbf{ANSI/TIA-568.2-D:} cableado de par trenzado balanceado.
    \item \textbf{ANSI/TIA-568.3-D:} cableado de fibra óptica.
    \item \textbf{ANSI/TIA-569-D:} rutas y espacios de telecomunicaciones.
    \item \textbf{ANSI/TIA-606-C:} administración y etiquetado del cableado.
    \item \textbf{ANSI/TIA-607-D:} puesta a tierra y enlace equipotencial para telecomunicaciones.
    \item \textbf{ISO/IEC 11801-1:} cableado genérico para instalaciones de clientes.
    \item \textbf{ISO/IEC 14763-2:} planificación, instalación y aseguramiento de calidad del cableado.
    \item \textbf{IEEE 802.3:} estándares de Ethernet para redes cableadas.
    \item \textbf{IEEE 802.11:} estándares de redes inalámbricas Wi-Fi.
    \item \textbf{Código Nacional de Electricidad (CNE) - Utilización (Perú):} criterios de seguridad eléctrica aplicables a canalizaciones y equipos.
    \item \textbf{Ley N.\textsuperscript{o} 29783 -- Ley de Seguridad y Salud en el Trabajo:} aplicable durante la ejecución de actividades de instalación y pruebas.
\end{itemize}

La aplicación de estas normas permite que la instalación de cableado y red se realice con criterios de seguridad, desempeño, interoperabilidad y calidad, garantizando condiciones adecuadas de conectividad para la facultad.
